\DocumentMetadata{
  pdfversion=2.0,
  pdfstandard=ua-2,
  lang=en-US,
  tagging=on,
  tagging-setup={math/setup=mathml-SE,tabsorder=structure}
}
\documentclass[11pt,letterpaper]{article}
\usepackage[margin=0.68in,includefoot]{geometry}
\usepackage{newcomputermodern}
\usepackage{amsmath,amscd,mathtools}
\providecommand{\square}{\mdlgwhtsquare}
\usepackage[hidelinks]{hyperref}
\hypersetup{
  pdftitle={Analysis PhD Exam, August 2012},
  pdfauthor={University of Florida Department of Mathematics},
  pdfsubject={Graduate mathematics examination},
  pdfdisplaydoctitle=true,
  bookmarksopen=true
}
\setcounter{secnumdepth}{0}
\setlength{\parindent}{0pt}
\setlength{\parskip}{0.28em}
\setlength{\emergencystretch}{3em}
\setlength{\leftmargini}{2.15em}
\setlength{\leftmarginii}{2.65em}
\setlength{\leftmarginiii}{3em}
\renewcommand{\labelenumi}{\textbf{\theenumi.}}
\pagestyle{empty}
\raggedbottom
\allowdisplaybreaks
\newenvironment{examproblems}[1]{%
  \begin{enumerate}%
  \setcounter{enumi}{#1}%
  \setlength{\topsep}{0.35em}%
  \setlength{\partopsep}{0pt}%
  \setlength{\itemsep}{0.48em}%
  \setlength{\parsep}{0.18em}%
}{\end{enumerate}}
\newenvironment{examsubparts}{%
  \begin{enumerate}%
  \setlength{\topsep}{0.18em}%
  \setlength{\partopsep}{0pt}%
  \setlength{\itemsep}{0.16em}%
  \setlength{\parsep}{0.1em}%
}{\end{enumerate}}
\newenvironment{examinstructions}{%
  \par\begingroup\itshape
}{%
  \par\endgroup\vspace{0.18em}
}
\makeatletter
\renewcommand\section{\@startsection{section}{1}{0pt}%
  {-1.25ex plus -.25ex minus -.1ex}%
  {0.55ex plus .1ex}%
  {\normalfont\large\bfseries}}
\renewcommand{\@maketitle}{%
  \newpage\null\vskip -1.2em%
  {\footnotesize\bfseries
    \href{https://www.ufl.edu/}{University of Florida}, \href{https://math.ufl.edu/}{Department of Mathematics}\par}%
  \vskip 0.18em%
  \tagstructbegin{tag=Title}%
    {\LARGE\bfseries\href{https://gma.math.ufl.edu/exams/analysis-phd/}{Analysis PhD Exam}, August 2012\par}%
  \tagstructend%
  \vskip 0.35em%
  \tagmcbegin{artifact}%
    \hrule height 1pt%
  \tagmcend%
  \vskip 0.55em%
}
\makeatother
\title{Analysis PhD Exam, August 2012}
\author{}
\date{}
\begin{document}
\maketitle
\thispagestyle{empty}
\begin{examinstructions}
Be sure to carefully present all work in a neat and logical fashion. Do not leave any gaps. State clearly theorems used in your proofs.
\end{examinstructions}
\begin{examproblems}{0}
\item
This problem has two parts.
\begin{examsubparts}
\item[\textnormal{(a)}]
State the Vitali convergence theorem (define all terms).
\item[\textnormal{(b)}]
Use this theorem to prove the Lebesgue Dominated Convergence theorem (first state this theorem).
\end{examsubparts}
\item
Let~\(f\) be an integrable function on \((S,\Sigma,\mu)\). Let~\(\nu\) be the indefinite integral of~\(f\). Let \(|\nu|(\cdot)\) denote the total variation of~\(\nu\) defined on~\(\Sigma\). Show that \(|\nu|(\cdot)\) is the indefinite integral of an a.e. unique integrable function~\(g\). What is the relationship between~\(g\) and~\(f\)? Prove.
\item
Let \((S,\Sigma,\mu)\) be a finite measure space. Suppose \(x^*\) belongs to the dual space of \(L^1(S,\Sigma,\mu)\). Show that there exist an a.e. unique integrable function~\(g\) such that 
\[
x^*(f)=\int_S fg\,d\mu,\quad\text{for each }f\in L^1(S,\Sigma,\mu).
\]
 [You need not prove that \(g\in L^\infty\), etc].
\item
Let \(\{y_n\}\) be an orthogonal sequence in a Hilbert space. Show \(\sum y_n\) converges unconditionally if and only if \(\sum |y_n|^2<\infty\).
\item
Let~\(X\) be a Banach space with closed linear subspaces~\(Y\) and~\(Z\). Suppose each~\(x\) in~\(X\) is the unique sum \(y+z\) where \(y\in Y\) and \(z\in Z\). Show that there is a constant~\(K\) such that \(|y|\le K|x|\) and \(|z|\le K|x|\) for each~\(x\) in~\(X\) with representation \(x=y+z\). [Hint: Use the open mapping theorem or closed graph theorem for certain maps]
\item
Let \(\{x_n\}\) be a weakly convergent sequence in a normed space~\(X\). Show that the weak limit belongs to the norm closed span of \(x_n\).
\end{examproblems}
\end{document}
